\scanpage{106}
\mspage{52}
11) Digression : déploiement d'un diviseur à croisements normaux.

Pour fixer les idées, je vais me placer dans le contexte des multiplicités analytiques lisses — mais on pourrait aussi, en prenant des multiplicités analytiques lisses \uline{relatives} — ou en remplaçant \unsure{l'hyp.} \unsure{analytiques|analytique} par des hyp. topologiques ad-hoc — ou en passant au contexte schématique\ldots.

Soit $X$ une multiplicité analytique lisse, et $D$ un diviseur sur $X$ à croisements normaux (hyp. purement locale en $(X,D)$, par définition). \deleted{Soit \unsure{$G$} un}

Pour tout pt $x$ de $D$, l'ens des ``branches'' de $D$ passant par $x$ est un ensemble fini $I(x)$, Par définition, les \uline{germes \inserted{en $x$} de sous-multiplicités} \uuline{\unsure{attachés|attachées}} \uuline{à $D$} sont les intersections d'un \unsure{nb|ens} fini \deleted{de} parmi les $D_{x,i}\in I(x)$.
\margin{il vaut mieux prendre un $x\in X$, et poser $I(x)=\varnothing$ si $x\notin D$}
Ces germes forment un ensemble ordonné \inserted{$\Delta_x\unsure{(=\Delta_x(X,D))}$} antiisomorphe à $\mathcal P(I(x))$, \deleted{Si \unsure{$J$ est}}\inserted{la codimension} \deleted{\unsure{une}} de $D_{x,J}$ $(J\subset I(x))$ étant donnée par
\[
(149)\qquad \operatorname{codim}(D_{x,J},X_x)=\operatorname{card}J
\]
\margin{105}
